\documentclass{article}
\usepackage{graphicx} % Required for inserting images
\usepackage{caption} % Required for \captionof
\usepackage{amsmath}
\usepackage{float}
\usepackage{comment}
\usepackage[T1]{fontenc}
\usepackage[utf8]{inputenc}
\usepackage{lmodern}
\usepackage{microtype}

\title{Hőtani mérés}
\author{Kern Luca, Csongor Vincze}
\date{2026 Április 1}

\begin{document}

\maketitle

\section*{4. Feladat: A termosz hőkapacitásának becslése}

%hibabecslés: nem a tömegmérés relatív hibája, hanem a 2-vel való osztásnak. Annak van egy 10-15-20%-os hibája. Ezt majd át kell gondolni!

Ebben a feladatban a termoszok hőkapacitását próbáltuk megállapítani. Ugyanis mint ahogy azt az előző feladatban is említettük, az hogy a termosz mennyi hőt vesz föl, egy olyan dolog, amely akár nagyban is befolyásolhatja, hogy hogyan alakulnak a hőcserék. Megmértük mindkét termosz tömegét, és felltételeztük, hogy a belső és a külső faluk ugyanolyan vastag,és mivel köztük szinte légüres tér van nincsen a belső és a külsőfalak között számottevő hőcsere. Ezért a termikus folyamatokban  résztvevő tömegeket a termoszok teljes tömegének felével közelítettük, amiket mérleggel mértünk meg. Ezután a hőkapacitást a $C=c*m$ összefüggés alapján számoltuk ki. A rozsdamentes acél fajhőjére az irodalmi adatot, azaz $c_{\text{rozsdamentes acél}}=500 \text{ J/(kg}^{\circ}\text{C)}$ használtuk föl.

\begin{table}[H]
    \centering
    \begin{tabular}{|l|c|c|}
        \hline
        Termosz & Teljes mért tömeg[g] & Becsült hőkapacitás [J/°C] \\\hline
        Zöld & 430 & 107,5 \\ \hline
        Fekete & 444 & 111 \\ \hline
        Piros & 380 & 95 \\ \hline
    \end{tabular}
    \caption{Termoszok hőkapacitásának becslése}
    \label{tab:hotan_4}
\end{table}

\begin{comment}
$m_{\text{fekete}}=444$ g\\
$m_{\text{zöld}}=430$ g\\
$m_{\text{piros}}=380$ g\\
$c_{\text{rozsdamentes acél}}=500 \text{ J/(kg}^{\circ}\text{C)}$\\
Becsült hőkapacitás:\\
$C_{\text{fekete}}=111$ J/°C\\
$C_{\text{zöld}}=107,5$ J/°C\\
$C_{\text{piros}}=95$ J/°C\\
\end{comment}

\section*{5. Feladat: A víz fajhőjének mérése elektromos fűtéssel}

%fontos: előtte és utána is mérni 1-2 percig!!!
Ebben feladatban a víz fajhőjét mértük ki, amit a következőképpen tettünk: a zöld termoszunkat megtöltöttük hideg vízzel, és a víz folyamatos melegítése mellett vizsgáltuk, hogy hogyan változik a hőmérésklet. A mérés kezdetén feljegyeztük a termoszban lévő fűtetlen víz hőméréskletét és tömegét. Ahhoz hogy a fűtőtestet megfelelő teljesítményen tudjuk üzemeltetni (50 W), megmértük az ellenállását. Majd 10 percig folyamatosan úgy melegítettük a termoszban lévő vizet, hogy kevertük az egyenletes elegyedés érdekében.
A mért hőmérséklet adatokat ábrázoltuk egy grafikonon az idő függvényében, amelyre egyenest illesztettünk, és meghatároztuk a meredekségét. A mért adataink alább találhatóak. Ezekből megállapítottuk a víz fajhőjét, amihez a termosz becsült hőkapacitását használtuk fel.

$V_{\text{víz}}=$0,77 L\\
$m_{\text{víz}}=$774 g\\
$R_{\text{fűtőtest}}=1,23 \Omega$\\
$U_{\text{fűtőtest}}=$7,95 V\\
$I_{\text{fűtőtest}}=6,46$ A\\
$T_{kezdeti}=26,4$ °C\\
$T_{veg}=36,4$ °C\\
termosz hőkapacitása az előző feladatból\\
$c_{\text{víz}}=4092 \text{ J/(kg}^{\circ}\text{C)}$\\

Ezek alapján a fajhő $c_{\text{víz}}=4532,52 \text{ J/(kg}^{\circ}\text{C)}$
-> Mi ez az elentmondás?????

Hibaszámítás: ..........



\begin{figure}[H]
    \centering
    \includegraphics[width=0.8\textwidth]{4_1_graf.pdf}
    \caption*{5. feladat: Fűtött víz hőmérsékletfüggése}
    \label{fig:5_1}
\end{figure}








\end{document}